\chapter{Teoretická východiska a analýza problematiky}
\label{sec:teorie}

Tato kapitola se věnuje teoretickým základům, na kterých je postaveno navržené testovací a monitorovací zařízení. Porozumění těmto konceptům je nezbytné pro správnou interpretaci dat získaných při analýze sítě.

\section{Protokoly síťové vrstvy a jejich výkonnostní rozdíly}
\label{sec:ipv4_ipv6}

Moderní počítačové sítě spoléhají na koexistenci protokolů IPv4 a IPv6 \cite{nastase2024}. Obě verze plní na síťové vrstvě tutéž úlohu, tedy adresování stanic a doručení paketu mezi sítěmi, liší se však natolik, že měřicí přístroj musí každou z nich obsluhovat samostatně. Souběžný provoz obou protokolů na jednom rozhraní, označovaný jako dvojí zásobník, je dnes v podnikových sítích běžný, a přenosný analyzátor proto musí umět ověřit funkčnost každé větve zvlášť.

Nejnápadnějším rozdílem je délka adresy, která vzrostla ze 32 na 128 bitů. Změna adresního prostoru ovšem není jediným zásahem do konstrukce protokolu. Hlavička protokolu IPv4 má proměnnou délku od 20 do 60 bajtů, neboť za povinnou částí mohou následovat volitelná pole, a obsahuje mimo jiné vlastní kontrolní součet a trojici polí sloužících k fragmentaci. Hlavička protokolu IPv6 je oproti tomu pevná a měří 40 bajtů \cite{rfc8200}. Veškerá doplňková funkcionalita byla přesunuta do rozšiřujících hlaviček, které se řetězí za základní hlavičku a na jejichž existenci upozorňuje pole Next Header.

Tento návrh má dva důsledky, jež se projeví právě při měření. Předně směrovač na cestě paketu vykonává méně práce: nepočítá znovu kontrolní součet, protože ten byl z hlavičky zcela odstraněn a spoléhá se na zabezpečení na linkové a transportní vrstvě, a nefragmentuje. Fragmentaci smí v protokolu IPv6 provést pouze odesílající stanice, která si k tomu musí předem zjistit nejmenší velikost přenosové jednotky na celé cestě. S tím souvisí požadavek, aby každá linka přenesla paket o velikosti alespoň 1280 bajtů. Druhým důsledkem je, že protokol IPv6 zcela opustil všesměrové vysílání a nahradil je vysíláním skupinovým, což se přímo promítá do způsobu, jakým zařízení v segmentu nacházejí jedno druhé; této otázce je věnována podkapitola \ref{sec:ipv6_adresace}.

Z pohledu měření propustnosti je podstatné, že dvojnásobná velikost základní hlavičky představuje režii, která se odečítá od užitečného zatížení. Při shodné velikosti rámce na linkové vrstvě zbývá na data o dvacet bajtů méně, a naměřená aplikační propustnost proto u protokolu IPv6 vychází nepatrně nižší. Jak dokládá srovnávací měření Tuifaigy a kolektivu \cite{tuifaiga2021}, protokol IPv4 v lokálních sítích skutečně dosahuje o něco vyšší propustnosti a nižší latence než IPv6, a to shodně v prostředí operačních systémů Windows i Linux.

Rozdíl je ovšem třeba interpretovat opatrně. Jeho relativní velikost závisí na tom, jak velké datové jednotky se přenášejí: u dlouhých toků s plnou velikostí rámce jde o jednotky procent, kdežto u provozu složeného z krátkých paketů může být znatelně vyšší. Naměřený rozdíl mezi oběma protokoly proto sám o sobě nevypovídá o kvalitě sítě, nýbrž o vlastnostech protokolu, a při vyhodnocování zátěžových testů popsaných v kapitole \ref{sec:realizace} je nutné jej od případných skutečných závad odlišit.

\section{Adresace a automatická konfigurace protokolu IPv6}
\label{sec:ipv6_adresace}

Rozdíl mezi oběma protokoly se neomezuje na velikost hlavičky. Podstatnější pro přenosný analyzátor je odlišný způsob, jakým se stanice ke své adrese vůbec dostane, neboť právě na něm závisí, zda přístroj dokáže oslovit protistranu bez jakékoli předchozí domluvy s její obsluhou. Této otázce je proto věnována samostatná podkapitola.

Adresa protokolu IPv6 má délku 128 bitů a její struktura je definována adresní architekturou \cite{rfc4291}, samotný protokol pak dokumentem \cite{rfc8200}. Zápis se člení do osmi skupin po čtyřech hexadecimálních číslicích, přičemž nejdelší souvislý úsek nulových skupin lze nahradit dvojicí dvojteček. Pro účely testovacího zařízení jsou podstatné tři kategorie adres. Adresa lokální linky z rozsahu \texttt{fe80::/10} vzniká na každém rozhraní automaticky ve chvíli, kdy je linka zdvižena, a platí pouze v rámci jednoho fyzického segmentu. Adresa globálního dosahu je směrovatelná v celém Internetu. Mezi nimi stojí kategorie jedinečných lokálních adres, která je pro navržené zařízení nejzajímavější.

Jedinečné lokální adresy jsou vymezeny prefixem \texttt{fc00::/7}, přičemž osmý bit rozlišuje způsob přidělení; při jeho nastavení, tedy pro prefix \texttt{fd00::/8}, si organizace generuje čtyřicetibitový identifikátor sama pseudonáhodným postupem \cite{rfc4193}. Takto vzniklý prefix je s velmi vysokou pravděpodobností jedinečný, není směrovatelný v Internetu a nezávisí na žádném poskytovateli připojení. Právě tyto vlastnosti odpovídají povaze přenosného analyzátoru, který pracuje v uzavřeném testovacím segmentu, nemá trvalé připojení k Internetu a jehož adresní plán nesmí kolidovat s adresami sítě, do níž je přístroj dočasně zapojen.

Druhou polovinu adresy tvoří identifikátor rozhraní o délce 64 bitů. Jeden ze způsobů jeho odvození, popsaný v adresní architektuře \cite{rfc4291}, vychází z fyzické adresy rozhraní a označuje se jako EUI-64. Osmačtyřicetibitová adresa se rozdělí na dvě poloviny, mezi něž se vloží dva bajty s hodnotou \texttt{fffe}, a v prvním bajtu se invertuje sedmý bit rozlišující univerzální a lokální správu. Fyzická adresa \texttt{00:e0:4c:68:02:16} tak vede na identifikátor \texttt{02e0:4cff:fe68:0216}. Tento postup je pro navržené zařízení prakticky významný: adresa protistrany z něj vychází deterministicky, a lze ji tedy dopočítat namísto toho, aby byla zapsána napevno do řídicí aplikace.

O zjišťování sousedů na společné lince se v protokolu IPv6 stará mechanismus Neighbor Discovery \cite{rfc4861}, který nahrazuje protokol ARP známý z protokolu IPv4 a je postaven nad protokolem ICMPv6. Zařízení si jeho prostřednictvím zjišťují vzájemné fyzické adresy, ověřují dosažitelnost sousedů a nacházejí směrovače v segmentu. Pro popisované zařízení jsou klíčové dvě z jeho zpráv. Zprávou Router Solicitation zasílanou na skupinovou adresu všech směrovačů se stanice po zdvižení rozhraní dotazuje, zda je v segmentu přítomen směrovač. Odpovědí je zpráva Router Advertisement, kterou směrovač inzeruje sám sebe a spolu s tím i prefixy platné v daném segmentu. Součástí mechanismu je rovněž detekce duplicitních adres, kterou stanice ověřuje, že jí zvolená adresa není v segmentu již obsazena.

Na zprávě Router Advertisement je postavena bezstavová automatická konfigurace adres \cite{rfc4862}. Stanice si z inzerovaného prefixu a z vlastního identifikátoru rozhraní sestaví úplnou adresu sama, bez jakéhokoli serveru, který by o přidělení vedl evidenci. Zda a jak tuto možnost využije, řídí trojice příznaků nesených ve zprávě, jejichž význam shrnuje Tabulka \ref{tab:priznaky}.

\begin{table}[htbp]
    \caption{Význam příznaků nesených ve zprávě Router Advertisement}
    \label{tab:priznaky}
    \centering
    \begin{tabular}{clp{7.5cm}}
        \hline
        Příznak & Název & Význam pro cílovou stanici \\
        \hline
        M & Managed address configuration & Adresu si stanice vyžádá od serveru DHCPv6 \\
        O & Other configuration & Ostatní údaje, například adresy serverů DNS, získá od serveru DHCPv6 \\
        A & Autonomous address configuration & Z tohoto konkrétního prefixu si stanice smí sestavit adresu sama \\
        \hline
    \end{tabular}
\end{table}

Zásadní je, že příznaky M a O se vztahují k celé zprávě, kdežto příznak A je nesen u každého inzerovaného prefixu zvlášť. Směrovač tedy může v jediné zprávě inzerovat několik prefixů a u každého z nich zvlášť rozhodnout, zda z něj stanice smí odvodit adresu vlastními silami. Obě cesty se přitom nevylučují a v praxi se běžně kombinují, takže jedna stanice může mít současně adresu sestavenou bezstavově a adresu přidělenou serverem.

Stavovou alternativou je protokol DHCPv6 \cite{rfc8415}, který přebírá roli známou z protokolu IPv4: server vede evidenci přidělených adres, váže je na identifikátor klienta a spravuje dobu jejich zapůjčení. Za cenu této evidence poskytuje správci úplnou kontrolu nad tím, která stanice jakou adresu obdrží, a umožňuje předat i další konfigurační údaje.

Rozdíl mezi oběma přístupy je pro přenosné testovací zařízení podstatný. Stanice s protokolem IPv4 bez běžícího klienta DHCP adresu nezíská nikdy, neboť jiný mechanismus protokol nenabízí. Naproti tomu jádro operačního systému s protokolem IPv6 přiděluje adresu lokální linky samočinně a při zdviženém příznaku A sestaví z inzerovaného prefixu i adresu širšího dosahu, aniž by na stanici musela běžet jakákoli služba pro konfiguraci sítě. Analyzátor schopný správně inzerovat prefix tak dokáže navázat spojení i s protistranou, na níž není nic nastaveno a nad níž nemá obsluha přístroje žádnou kontrolu. Ověření tohoto předpokladu na sestaveném prototypu je popsáno v podkapitole \ref{sec:overeni_autokonfigurace}.

\section{Monitorování, generování a bezpečnostní audit síťového provozu}
\label{sec:monitorovani}

Práci se síťovým provozem lze rozdělit na dva zásadně odlišné přístupy. Pasivní monitorování provoz pouze pozoruje a do sítě nic nevysílá, kdežto aktivní měření provoz samo vytváří a sleduje, jak se s ním síť vypořádá. Přenosný analyzátor musí zvládnout oba, neboť odpovídají na jiné otázky: pasivní sledování ukáže, co se v segmentu skutečně děje, aktivní měření zjistí, čeho je segment schopen.

Pasivní zachytávání je v systémech na bázi Linuxu postaveno na knihovně libpcap a na filtrovacím mechanismu zabudovaném v jádře, který umožňuje odmítnout nezajímavé pakety ještě dříve, než jsou předány aplikaci. Tím se výrazně snižuje zatížení procesoru, což je na jednodeskovém počítači s omezeným výkonem podstatné. Aby rozhraní přijímalo i rámce, které nejsou určeny jemu, musí být přepnuto do promiskuitního režimu. Nad touto vrstvou pracuje nástroj tcpdump, který zachycený provoz vypisuje nebo ukládá do souboru ve formátu pcap pro pozdější rozbor. Zásadním omezením pasivního odposlechu je topologie moderních sítí: přepínač doručuje rámce pouze na port, kam náleží cílová adresa, takže analyzátor zachytí buď provoz určený přímo jemu, nebo provoz na zrcadleném portu, případně na lince, do níž je vložen přímo. Tuto skutečnost je nutné mít na paměti při výkladu naměřených statistik.

Pro úlohy, kde nestačí provoz jen pozorovat, ale je třeba pakety sestavovat na míru, slouží knihovna Scapy v jazyce Python. Umožňuje poskládat rámec vrstvu po vrstvě, odeslat jej a zpracovat odpovědi, čehož využívá i mapování segmentu protokolem ARP: analyzátor rozešle dotazy na všechny adresy zvoleného rozsahu a ze získaných odpovědí sestaví seznam aktivních stanic i s jejich fyzickými adresami. Jak uvádí Carter \cite{carter2024}, právě automatizace těchto úloh pomocí jazyka Python zásadně rozšiřuje možnosti síťového auditu, neboť dovoluje vlastní postupy přizpůsobit konkrétnímu prostředí.

Aktivní stranu měření zastupuje nástroj iperf3, pracující v uspořádání klient a server. Volba transportního protokolu určuje, co vlastně měříme. Měření nad protokolem TCP zjišťuje propustnost dosažitelnou pro aplikaci včetně vlivu řízení toku a zotavení ze ztrát, takže výsledek odráží chování reálného přenosu dat. Měření nad protokolem UDP naopak vysílá předem zvolenou rychlostí bez ohledu na stav sítě a vedle propustnosti vypovídá o ztrátovosti a rozptylu zpoždění. Rozdílem mezi spojovanou a nespojovanou komunikací při vysokých rychlostech se podrobně zabývají Gál a kolektiv \cite{gal2021}, jejichž práce tvoří metodický základ pro měření provedená v této práci. Pro vypovídající výsledek je nutné dodržet dostatečnou dobu trvání testu, aby se řízení toku stačilo ustálit, a měření opakovat.

Třetí skupinu tvoří nástroje pro aktivní bezpečnostní průzkum, mezi nimiž má výsadní postavení mapovač Nmap. Jeho práce probíhá v několika stupních. Nejprve zjistí, které adresy ve zvoleném rozsahu odpovídají, poté u nalezených stanic zkoumá stav vybraných portů, a nakonec se z podoby odpovědí pokusí odvodit, jaké služby na nich naslouchají a jaký operační systém stanice provozuje. Poslední jmenovaná úloha, označovaná jako otisk zásobníku, vychází z drobných odlišností v tom, jak jednotlivé implementace sestavují odpovědi na neobvyklé kombinace příznaků.

S aktivním průzkumem je nerozlučně spjata otázka oprávněnosti. Skenování cizí sítě bez souhlasu jejího správce je nepřípustné, a to bez ohledu na to, zda způsobí škodu. Navržené zařízení tuto otázku řeší svou konstrukcí: pracuje v uzavřeném testovacím segmentu, který sám vytváří a jehož jediným dalším účastníkem je stanice připojená přímo kabelem. Skenování tak nikdy neopustí prostor vymezený obsluhou přístroje, což je vedle přesnosti měření druhým důvodem pro logické oddělení popsané v následující podkapitole.

\section{Architektura Linux Network Namespaces (netns)}
\label{sec:netns_teorie}

Pro realizaci autonomního testování na jediném fyzickém stroji se využívá technologie Linux Network Namespaces (netns). Tato technologie poskytuje logickou izolaci síťových prostředků přímo na úrovni jádra OS. Každý jmenný prostor udržuje své vlastní nezávislé instance síťového zásobníku, což zahrnuje izolované směrovací tabulky, pravidla brány firewall a nezávislá síťová rozhraní.

Při vytvoření nového jmenného prostoru je tento prostor zpočátku zcela prázdný a obsahuje pouze rozhraní zpětné smyčky, které je navíc ve vypnutém stavu. Teprve přesunutím fyzického rozhraní do tohoto prostoru a jeho zdvižením vznikne funkční síťové prostředí. Přesunuté rozhraní zároveň zmizí z původního prostoru, takže hostitelský systém o něm nadále neví a nemůže jeho prostřednictvím nic odeslat.

Izolace se týká celého síťového zásobníku, nikoli pouze rozhraní. Každý prostor má vlastní směrovací tabulky, vlastní tabulku sousedů, vlastní sadu pravidel paketového filtru i vlastní naslouchající sokety. Dva procesy běžící v odlišných prostorech tedy mohou obsadit tentýž port, aniž by si navzájem překážely, což je pro měřicí přístroj s několika instancemi téže služby výhodné. Příkazy se do zvoleného prostoru směrují pomocnou utilitou, která proces spustí přímo v něm; celý mechanismus je součástí standardního jádra a nevyžaduje žádnou nadstavbu.

Mají-li spolu jmenné prostory naopak komunikovat, propojí se dvojicí virtuálních rozhraní typu veth, která se chová jako pomyslný kabel mezi nimi. V případě popisovaného analyzátoru je však žádoucí pravý opak, a proto žádné takové propojení vytvořeno není. Důsledkem je, že mezi oběma testovacími porty neexistuje cesta a provoz z jednoho se do druhého nedostane ani omylem.

Tutéž úlohu by bylo možné řešit i jinak. Nabízí se použití dvou samostatných počítačů, což ovšem odporuje požadavku na přenositelnost a zvyšuje cenu. Virtualizace by přinesla režii, kterou jednodeskový počítač nese jen obtížně. Zbývá cesta pravidel pro směrování podle zdrojové adresy nebo svázání aplikace s konkrétním rozhraním; tato řešení sice fungují, ale izolace je u nich věcí správné konfigurace a jediná chyba či nástroj, který svázání s rozhraním nepodporuje, ji zruší. Jmenné prostory oproti tomu vynucují oddělení na úrovni jádra, takže nesprávně nasměrovaný provoz nemá kudy odejít a chyba se projeví okamžitě a jednoznačně, nikoli tichým zkreslením naměřených hodnot.

\section{Analýza problematiky a existujících řešení}
\label{sec:analyza}

Úloha ověřit stav síťové zásuvky, změřit dosažitelnou propustnost nebo zjistit, jaká zařízení se v segmentu nacházejí, se v praxi řeší třemi odlišnými cestami. Každá z nich má své přednosti a právě z jejich porovnání vyplývá zadání této práce.

Nejrozšířenější cestou je univerzální počítač vybavený otevřenými nástroji. Zachytávání provozu obstará tcpdump nebo Wireshark, zátěžové testy iperf3 a průzkum segmentu Nmap. Toto řešení nabízí největší výkon i volnost, protože měřicí stanicí je plnohodnotný počítač. Jeho slabinou je ovšem provozní stránka věci. Notebook je v rozvaděči nebo pod stolem těžko použitelný, k propojení s testovanou sítí obvykle potřebuje redukci a před každým měřením je nutné ručně nastavit adresy a ověřit, kudy provoz skutečně poteče. Připojí-li se navíc technik do téhož segmentu dvěma rozhraními, aby mohl současně naslouchat a generovat zátěž, naráží na konflikt ve směrovacích tabulkách popsaný v podkapitole \ref{sec:netns_teorie}: operační systém odešle provoz tím rozhraním, které odpovídá výchozí trase, nikoli tím, které měl technik na mysli.

Druhou cestou jsou komerční příruční testery určené právě pro práci v terénu. Jsou ergonomické, po zapnutí okamžitě připravené a jejich obsluha nevyžaduje znalost příkazové řádky. Daní za tuto pohodlnost je uzavřenost. Rozsah funkcí je pevně dán výrobcem, vlastní skript nebo nový testovací postup do přístroje doplnit nelze a rozšíření obvykle znamená nákup vyššího modelu. Pro výukové a laboratorní prostředí, kde je žádoucí sledovat i to, jak měření vlastně probíhá, je uzavřená konstrukce spíše překážkou.

Třetí cestou, které se věnuje i tato práce, je jednodeskový počítač s otevřeným operačním systémem. Vhodnost platformy Raspberry Pi pro tento účel je v odborné literatuře doložena. Coşar a Kiran \cite{cosar2018} se zaměřili přímo na výkonnostní stránku věci a měřili chování Raspberry Pi při analýze síťového provozu včetně zatížení procesoru a operační paměti. Tripathi a Kumar \cite{tripathi2018} na téže platformě provozovali současně systém pro detekci průniků, past na útočníky a analyzátor paketů, čímž ukázali, že výpočetní výkon desky postačuje i pro souběh několika síťových služeb. Hariawan a Sunaringtyas \cite{hariawan2021} obdobné řešení postavili na novější generaci desky a nasadili je v prostředí domácí sítě.

Uvedené práce mají společný rys. Do popředí staví pasivní stranu úlohy, tedy zachytávání, detekci a vyhodnocování provozu, který v síti vzniká nezávisle na měřicím zařízení. Aktivní generování zátěže a zejména jeho oddělení od měřicí cesty v nich ústřední roli nehraje, stejně jako otázka, jak přístroj ovládat bez připojeného počítače.

Právě zde se otevírá prostor pro zadání této práce. Cílem je spojit přednosti obou krajních přístupů: zachovat otevřenost a rozšiřitelnost počítače s operačním systémem Linux a zároveň dosáhnout ergonomie příručního testeru. Z porovnání existujících řešení plynou čtyři požadavky na navrhované zařízení. Přístroj musí umožnit současné pasivní naslouchání a aktivní generování provozu, přičemž obě činnosti musí být vedeny odděleně, aby generovaná zátěž nezkreslovala zachycené statistiky. Toto oddělení musí být logicky vynuceno, nikoli ponecháno na kázni obsluhy, k čemuž slouží jmenné prostory. Ovládání musí být možné přímo na zařízení, bez připojeného notebooku a bez fyzické klávesnice. A konečně musí být protokol IPv6 podporován ve stejné hloubce jako IPv4, včetně automatické konfigurace adres na protistraně, neboť právě ta rozhoduje o tom, zda přístroj dokáže oslovit stanici, na níž není nic připraveno.

Naplnění těchto požadavků je předmětem následující kapitoly.